%% 6-5 空気抵抗と終端速度（第14回：組 V6-3。ページ制限なし。偶数ページ）
%% ほかの巻が指している形：速さに比例する抵抗 f=bv、2乗に比例する抵抗 f=cv^2（b、c は抵抗の係数）、下向きが＋、
%%   m dv/dt = mg - bv、終端速度 v_f = mg/b、v = v_f(1-e^{-t/τ})、τ = m/b、2乗のときの終端速度 v_t = √(mg/c)（物理 p1-02）、
%%   スカイダイバー 80 kg・c=0.26・v_t≈54.9 m/s（物理 p1-02・6-11）、エレベーター（5-5：0〜2秒 +1 m/s²、2〜8秒 0、8〜10秒 −1 m/s²）、
%%   RL 回路の最後の電流 E/R ＝「電流の終端速度」（6-4 の RL と同じ数：12 V・400 Ω・0.2 H）。
\makeatletter
% 本文の小見出し（太字の1行）
\newcommand\bkvifhead[1]{\par\addvspace{6pt}\Needspace{6\baselineskip}\noindent{\sffamily\bfseries\color{\bkTT}#1}\par\nobreak\vspace{2pt}}
% 大事な式のわく
\newcommand\bkvifframe[1]{\par\addvspace{4pt}{\centering\tikz\node[draw=\bkTT, line width=0.9pt, fill=bkTheme!10, rounded corners=1.5mm, inner xsep=6mm, inner ysep=2.5mm, align=center]{#1};\par}\addvspace{4pt}\noindent\ignorespaces}
% 例題の途中式の行
\newcommand\bkvifq{\par\vspace{1pt}\noindent\hangindent=2.6\zw\hangafter=0\relax}
% 雨つぶ  \bkvifdrop{(x,y)}{大きさ}
\newcommand\bkvifdrop[2]{\begin{scope}[shift={#1}, scale=#2]
  \filldraw[fill=bkBlue!35, draw=bkBlue!80!black, line width=0.6pt] (0,2.4) .. controls (0.6,1.2) and (1.3,0.6) .. (1.3,-0.2) arc (0:-180:1.3) .. controls (-1.3,0.6) and (-0.6,1.2) .. (0,2.4) -- cycle;
\end{scope}}
% 雲  \bkvifcloud{(x,y)}
\newcommand\bkvifcloud[1]{\begin{scope}[shift={#1}]
  \fill[bkGray!25] (0,0) circle (4) (5,1.5) circle (5) (10.5,0.5) circle (4.2) (14.5,-0.5) circle (3) (5,-2) ellipse (9 and 2.6);
\end{scope}}
% コーヒーフィルター（横から見た形）  \bkviffilter{(x,y)}{枚数}
\newcommand\bkviffilter[2]{\begin{scope}[shift={#1}]
  \foreach \k in {1,...,#2} {\filldraw[fill=white, draw=bkInk!70, line width=0.5pt] ({-5},{0.7*(\k-1)}) -- (5,{0.7*(\k-1)}) -- (3.2,{0.7*(\k-1)-4}) -- (-3.2,{0.7*(\k-1)-4}) -- cycle;}
  \foreach \x in {-2.4,-0.8,0.8,2.4} {\draw[bkInk!35, line width=0.3pt] ({\x*1.25},{0.7*(#2-1)}) -- (\x,{0.7*(#2-1)-4});}
\end{scope}}
% 回路の部品（6-4 と同じ形）
\newcommand\bkvifR[2]{\begin{scope}[shift={#1}, rotate=#2]
  \filldraw[fill=white, draw=bkInk, line width=0.7pt] (-4.5,-1.6) rectangle (4.5,1.6);\end{scope}}
\newcommand\bkvifL[2]{\begin{scope}[shift={#1}, rotate=#2]
  \fill[white] (-5.2,-0.6) rectangle (5.2,3.0);
  \draw[bkInk, line width=0.7pt] (-5,0) arc (180:0:1.25) arc (180:0:1.25) arc (180:0:1.25) arc (180:0:1.25);\end{scope}}
\newcommand\bkvifE[2]{\begin{scope}[shift={#1}, rotate=#2]
  \fill[white] (-1.2,-4.2) rectangle (1.2,4.2);
  \draw[bkInk, line width=0.7pt] (-1.0,-4.0) -- (-1.0,4.0);
  \draw[bkInk, line width=1.8pt] (1.0,-2.0) -- (1.0,2.0);\end{scope}}
\newcommand\bkvifS[2]{\begin{scope}[shift={#1}, rotate=#2]
  \fill[white] (-3.5,-0.6) rectangle (3.5,3.6);
  \fill[bkInk] (-3,0) circle (0.6) (3,0) circle (0.6);
  \draw[bkInk, line width=0.8pt] (-3,0) -- (2.6,3.0);\end{scope}}
\makeatother

\bkchapter{6}{5}{空気抵抗と終端速度}{\bklevelrange{2}{3}}{空気抵抗と終端速度}
\bkmichishirube{5-5,5-13,6-2,6-4}{6-11,6-15}

\begin{goal}
\begin{itemize}[label={\color{\bkTT}●}, leftmargin=1.6\zw, itemsep=0pt, topsep=1pt]
\item \bklevel{2}\enspace 落ちる物にはたらく力から式を立て、終端速度を「速さが変わらなくなる」ことから求められる。
\item \bklevel{2}\enspace 表計算で、空気の抵抗を受けて落ちる物の速さを計算できる。
\item \bklevel{3}\enspace 速さに比例する抵抗の式を解いて $v=v_{\mathrm f}(1-e^{-t/\tau})$ を作り、エレベーターや RL 回路にも使える。
\end{itemize}
\end{goal}

\begin{hatena}
\noindent\begin{minipage}[c]{0.62\linewidth}\raggedright
雨つぶは、高さ1000\,mほどの雲から落ちてきます。空気がなければ、地面に着くときの速さは秒速140\,m（時速504\,km）。新幹線よりずっと速く、当たったら大けがをしそうです。

でも、本当の雨つぶ（直径2\,mm）は、秒速6.5\,mほどで落ちてきます。雨つぶは、なぜどこまでも速くならないのでしょう。
\end{minipage}\hfill
\begin{minipage}[c]{0.35\linewidth}\centering
\begin{tikzpicture}[x=1mm, y=1mm]
  \bkvifcloud{(4,40)}
  \bkvifdrop{(10,30)}{1}
  \bkvifdrop{(10,20)}{1}
  \bkvifdrop{(10,10)}{1}
  \draw[bkGray, line width=0.8pt] (-2,0) -- (30,0);
  \draw[bkInk!60, line width=0.6pt, {Stealth[length=3pt]}-{Stealth[length=3pt]}] (24,1) -- (24,35.5);
  \node[font=\sffamily\footnotesize, anchor=west, fill=bkYellowBg, inner sep=1pt] at (25,18) {1000\,m};
  \node[font=\sffamily\footnotesize, anchor=west] at (13,10) {？};
\end{tikzpicture}
\end{minipage}
\end{hatena}

\noindent 5-13では、空気がないとして落ちる物を調べ、$v=gt$（速さは時間に比例してふえる）を求めました。この章では、空気の抵抗を入れて、落ちる物の式を立てます。速さが変わらなくなる「終端速度」は、式を解かなくても求められます。そして、6-4の「目標の値へ近づく形」を使って、速さの変わり方を式で求めます。

式が少し長くなりますが、1回でわからなくても大丈夫です。まずコーヒーフィルターを落として、目で確かめてから式に進みます。% Ver2 No.24 声かけ

%% ---------------- やってみよう・思い出そう ----------------
\begin{yattemiyou}[コーヒーフィルターを落とそう]
\textbf{用意するもの}\par
紙のコーヒーフィルター（同じもの）7枚、メジャー、スマートフォン（動画をとる）
\begin{enumerate}[label={\bkcircled[bkGreen]{\arabic*}}, leftmargin=2.4\zw, itemsep=1pt, topsep=3pt]
\item かべに、床から1.5\,mの所に目じるしをつけよう。
\item 2.5\,mほどの高さから、フィルター1枚を口を上にして落とし、動画をとろう。
\item 動画をコマ送りして、目じるしから床までの時間を読もう。
\item 2枚、4枚を重ねて、同じようにはかろう。速さ＝1.5\,m÷時間を計算しよう。
\end{enumerate}
\noindent\begin{minipage}[c]{0.32\linewidth}\centering
\begin{tikzpicture}[x=1.12mm, y=1.12mm]
  \bkviffilter{(5,14)}{1}
  \bkviffilter{(17,14)}{2}
  \bkviffilter{(29,14)}{4}
  \node[font=\sffamily\footnotesize, anchor=north] at (5,8.5) {1枚};
  \node[font=\sffamily\footnotesize, anchor=north] at (17,8.5) {2枚};
  \node[font=\sffamily\footnotesize, anchor=north] at (29,8.5) {4枚};
\end{tikzpicture}
\end{minipage}\hfill
\begin{minipage}[c]{0.66\linewidth}\centering
\begin{bktable}{colspec={Q[l,m,wd=22mm] *{3}{Q[c,m,wd=11mm]}}, rowsep=1pt, colsep=3pt}
重ねた枚数 & 1 & 2 & 4 \\
時間（s） & 1.42 & 1.02 & 0.73 \\
速さ（m/s） & 1.06 & 1.47 & 2.05 \\
\end{bktable}\par\vspace{2pt}
{\sffamily\footnotesize ミナさんがはかった結果}
\end{minipage}
\end{yattemiyou}

\noindent フィルターは、手をはなしてすぐに一定の速さになり、目じるしから下はほぼ同じ速さで落ちていきました。重ねて重くすると、速く落ちます。でも、重さを2倍にしても速さは2倍にならず、約1.4倍（$1.47\div1.06=1.39$）です。4倍にして、やっと約2倍（$2.05\div1.06=1.93$）。この「1.4」と「2」はどこから来るのでしょう。
\begin{bkfigure}
\begin{tikzpicture}[x=1mm, y=1mm]
  \foreach \t/\a/\b in {0/0/0, 2/0.135/0.172, 3/0.239/0.341, 4/0.345/0.532, 5/0.451/0.731, 6/0.557/0.934} {
    {\bkviffilter{(25,{-70*\a+2})}{1}}
    {\bkviffilter{(75,{-70*\b+2})}{4}}
    \node[font=\sffamily\footnotesize, anchor=east] at (17,{-70*\a}) {\ifnum\t=0 0\else 0.\t\fi\,s};
    \node[font=\sffamily\footnotesize, anchor=west] at (83,{-70*\b}) {\ifnum\t=0 0\else 0.\t\fi\,s};}
  \node[font=\sffamily\footnotesize, anchor=south] at (25,5.5) {1枚};
  \node[font=\sffamily\footnotesize, anchor=south] at (75,5.5) {4枚};
  \draw[bkRed, line width=0.8pt, {Stealth[length=3pt]}-{Stealth[length=3pt]}] (34,{-70*0.451}) -- (34,{-70*0.557});
  \node[font=\sffamily\footnotesize, text=bkRed, anchor=west] at (35,{-70*0.504}) {0.1\,sに約11\,cm};
  \draw[bkRed, line width=0.8pt, {Stealth[length=3pt]}-{Stealth[length=3pt]}] (63,{-70*0.731}) -- (63,{-70*0.934});
  \node[font=\sffamily\footnotesize, text=bkRed, anchor=east] at (62,{-70*0.832}) {0.1\,sに約20\,cm};
\end{tikzpicture}
\bkfigcaption{0.1秒ごとのフィルターの位置（計算した図）。はじめの0.2秒ほどで間かくがそろい、そのあとは同じ速さで落ちる。4枚は1枚の約2倍の間かく。}
\end{bkfigure}

\begin{omoidasou}
\bkrecall{5-5}{運動の法則 $F=ma$：物にはたらく力の合計＝質量×加速度。加速度は速度の変わる速さ $a=\dfrac{\dd v}{\dd t}$。}
\bkrecall{5-13}{空気がなければ、どんな物も同じ加速度 $g=\qty{9.8}{m/s^2}$ で落ちる：$v=gt$、$x=\frac{1}{2}gt^2$、高さ $h$ 落ちたときの速さ $\sqrt{2gh}$。}
\bkrecall{6-2}{オイラー法：次の速さ＝今の速さ＋（今の加速度）×$\Delta t$。}
\bkrecall{6-4}{$\dfrac{\dd y}{\dd t}=\dfrac{y_\infty-y}{\tau}$ の答えは $y=y_\infty+(y_0-y_\infty)e^{-t/\tau}$。最後の値は $\dfrac{\dd y}{\dd t}=0$ とおけばわかる。}
\end{omoidasou}

%% ---------------- 数学の道具 ----------------
\bksection{doguu}{数学の道具}

\bksubsection{1}{落ちる物にはたらく力と、運動の式}

\leavevmode\label{bk:6-5:shiki}空気の中を落ちる雨つぶには、2つの力がはたらいています。
\begin{itemize}[leftmargin=1.2\zw, itemsep=1pt, topsep=2pt, label={\color{\bkTT}●}]
\item \textbf{重力}（下向き）：大きさ $mg$。$m$〔kg〕は質量、$g=\qty{9.8}{m/s^2}$ は地球の上ではかった重力の加速度です。質量は「物の量」で、どこでも変わりません。重力 $mg$〔N〕は地球が物を引く力で、これを「重さ」ともいいます。
\item \textbf{空気の抵抗}（上向き）：大きさを $f$〔N〕とします。落ちる物は空気をおしのけながら進むので、空気からおし返されます。抵抗の向きは、いつも動きと反対です。
\end{itemize}
\begin{bkfigure}
\begin{tikzpicture}[x=1mm, y=1mm]
  \foreach \px/\f/\ttl/\v in {0/0/{落ち始め}/{$v=0$}, 42/5/{とちゅう}/{$v$ が小さい}, 84/10/{終端速度}/{$v=v_{\mathrm f}$}} {
    \begin{scope}[xshift=\px mm]
      \bkvifdrop{(14,14)}{2.4}
      \draw[bkRed, line width=1.6pt, -{Stealth[length=5pt]}] (14,8) -- (14,-2);
      \node[font=\sffamily\footnotesize, text=bkRed, anchor=west] at (15.5,0) {$mg$};
      \ifnum\f>0 \draw[bkBlue, line width=1.6pt, -{Stealth[length=5pt]}] (14,21) -- (14,{21+\f});\fi
      \node[font=\sffamily\footnotesize, anchor=north] at (14,-4) {\ttl};
      \node[font=\sffamily\footnotesize, anchor=north] at (14,-9) {\v};
    \end{scope}}
  \node[font=\sffamily\footnotesize, text=bkBlue, anchor=west] at (57.5,25) {$f$（小さい）};
  \node[font=\sffamily\footnotesize, text=bkBlue, anchor=west] at (99.5,27) {$f=mg$};
  \node[font=\sffamily\footnotesize, anchor=west] at (15.5,25) {抵抗0};
  \draw[bkInk!50, line width=0.6pt, -{Stealth[length=4pt]}] (118,22) -- (118,4);
  \node[font=\sffamily\footnotesize, anchor=north] at (118,3) {下が$+$};
\end{tikzpicture}
\bkfigcaption{落ちる雨つぶにはたらく力。重力 $mg$（赤）はいつも同じ。抵抗 $f$（青）は速くなるほど大きくなり、やがて重力と同じ大きさになる。}\label{bk:6-5:fig-chikara}
\end{bkfigure}
下向きを $+$ にします（物理シリーズ『力学』と同じ）。下向きの力の合計は $mg-f$ です。5-5の運動の法則「力の合計＝質量×加速度」に入れると、
\bkvifframe{$m\,\dfrac{\dd v}{\dd t}=mg-f$}
これが、空気の中を落ちる物の式です。\bkidx{うんどうほうていしき}{運動方程式}$\dfrac{\dd v}{\dd t}$ が入っているので、6-1で学んだ微分方程式です。空気がなければ $f=0$ で、$\dfrac{\dd v}{\dd t}=g$。5-13の自由落下にもどります。

\bksubsection{2}{抵抗の大きさは、速さで変わる}

\leavevmode\label{bk:6-5:teikou}抵抗は、止まっていれば0で、速くなるほど大きくなります。自転車で走ると、速いほど強い向かい風を感じるのと同じです。抵抗の大きさは、主に次の2つの形で表します。
\begin{bkcenter}
\begin{bktable}{colspec={Q[l,m,wd=24mm] Q[c,m,wd=22mm] X[l,m]}, width=\linewidth, rowsep=2pt}
形 & 式 & どんなとき \\
速さに比例 & $f=bv$ & ごく小さい物がゆっくり動くとき（霧のつぶ、ほこり、ねばりけの強い油の中の小さな玉） \\
速さの2乗に比例 & $f=cv^2$ & ふつうの大きさの物が速く動くとき（雨つぶ、ボール、自動車、人） \\
\end{bktable}
\end{bkcenter}
$b$、$c$ を\textbf{抵抗の係数}といいます。\bkidx{ていこうのけいすう}{抵抗の係数}どちらになるか、いくつになるかは、物の大きさ・形・速さ、空気のねばりけで決まります（物理シリーズ『力学』）。文字を $b$、$c$ にしたのは、ばね定数の $k$ とまぎれないようにするためです（物理シリーズと同じ文字）。
\begin{bkfigure}
\begin{tikzpicture}
\begin{axis}[bkplot, width=86mm, height=46mm, xmin=0, xmax=3.4, ymin=0, ymax=10.5, xtick={1,2,3}, ytick=\empty, xlabel={速さ $v$}, ylabel={抵抗 $f$}]
  \addplot[bkcurve, domain=0:3.2] {2*x};
  \addplot[bkcurve2, domain=0:3.2] {x^2};
  \node[font=\sffamily\footnotesize, text=bkBlue, anchor=west] at (axis cs:3.22,6.4) {$f=bv$};
  \node[font=\sffamily\footnotesize, text=bkOrange!80!black, anchor=west] at (axis cs:3.22,10.2) {$f=cv^2$};
  \addplot[only marks, mark=*, mark size=1.5pt, bkInk] coordinates {(1,1) (2,4) (3,9)};
  \node[font=\sffamily\footnotesize, anchor=south east] at (axis cs:2,4.4) {2倍で4倍};
  \node[font=\sffamily\footnotesize, anchor=south east] at (axis cs:3,9.4) {3倍で9倍};
\end{axis}
\end{tikzpicture}
\bkfigcaption{速さに比例する抵抗（青）は、速さが2倍で2倍。2乗に比例する抵抗（オレンジ）は、速さが2倍で4倍、3倍で9倍。}
\end{bkfigure}

\bkvifhead{なぜ2乗になるの？}
速く進む物は、2つの理由で、空気から強くおし返されます。
\begin{enumerate}[label={\bkcircled[\bkTF]{\arabic*}}, leftmargin=2.4\zw, itemsep=1pt, topsep=2pt]
\item 1秒間におしのける空気の量が、速さに比例してふえる（2倍の速さなら、2倍の空気をおしのける）。
\item 1つ1つの空気を、速さに比例して強くはじく（2倍の速さなら、2倍の勢いではじく）。
\end{enumerate}
2つが重なって、$2\times2=4$ 倍。だから、速さの2乗に比例します。小さい物がゆっくり動くときは、空気のねばりけでまわりの空気を引きずる力のほうが大きく、こちらは速さに比例します。

\bkvifhead{単位を確かめる}
力 $f$ の単位はN（ニュートン）です。$f=bv$ なら、$b=\dfrac{f}{v}$ の単位は $\dfrac{\text{N}}{\text{m/s}}=\text{N·s/m}$。$f=cv^2$ なら、$c=\dfrac{f}{v^2}$ の単位は $\dfrac{\text{N}}{(\text{m/s})^2}=\text{N·s}^2\text{/m}^2$。$b$ と $c$ は、単位のちがう別の量です。

\bksubsection{3}{解かなくてもわかる：終端速度}

\leavevmode\label{bk:6-5:shuutan}図\ref{bk:6-5:fig-chikara}のように、落ち始めは速さが0なので、抵抗も0。加速度は $g$ です。速くなると抵抗がふえ、下向きの力の合計 $mg-f$ は小さくなるので、速くなり方がにぶくなります。やがて抵抗が重力と同じ大きさになると、力の合計が0になり、速さはもう変わりません。この速さを\textbf{終端速度}（しゅうたんそくど）といいます。\bkidx{しゅうたんそくど}{終端速度}

\noindent 速さが変わらないので、$\dfrac{\dd v}{\dd t}=0$ です。6-4の「最後の値」と同じ考えで、式の左辺を0とおきます。
\begin{bkhenkei}
0=mg-bv_{\mathrm f} & 速さに比例するとき：$\frac{\dd v}{\dd t}=0$ とおく \\
bv_{\mathrm f}=mg & $bv_{\mathrm f}$ を左辺へ移し、左右を入れかえる \\
v_{\mathrm f}=\frac{mg}{b} & 両辺を $b$ でわる \\
0=mg-cv_{\mathrm t}^2 & 2乗に比例するとき：$\frac{\dd v}{\dd t}=0$ とおく \\
cv_{\mathrm t}^2=mg & $cv_{\mathrm t}^2$ を左辺へ移し、左右を入れかえる \\
v_{\mathrm t}^2=\frac{mg}{c} & 両辺を $c$ でわる \\
v_{\mathrm t}=\sqrt{\frac{mg}{c}} & 正のほうの平方根をとる（3-3） \\
\end{bkhenkei}
2つの終端速度を区別するため、速さに比例するときは $v_{\mathrm f}$、2乗に比例するときは $v_{\mathrm t}$ と書きます（物理シリーズと同じ書き方）。微分方程式を解かなくても、「変わらなくなる速さ」だけなら、かけ算・わり算と平方根で求められます。
\bkvifframe{速さに比例\quad $v_{\mathrm f}=\dfrac{mg}{b}$\qquad 速さの2乗に比例\quad $v_{\mathrm t}=\sqrt{\dfrac{mg}{c}}$}

\bkvifhead{\bkhatenaref の答え}
雨つぶは、速くなるほど空気の抵抗が大きくなり、抵抗が重力とつり合う終端速度に達すると、もう速くなりません。直径2\,mmの雨つぶでは、質量 $m=4.2\times10^{-6}$\,kg（0.0000042\,kg）、抵抗の係数 $c\approx9.7\times10^{-7}$\,N·s$^2$/m$^2$（形と大きさから計算した値）です。$10^{-6}\div10^{-7}=10$ なので、
\[ v_{\mathrm t}=\sqrt{\frac{4.2\times10^{-6}\times9.8}{9.7\times10^{-7}}}=\sqrt{\frac{4.2\times9.8\times10}{9.7}}=\sqrt{42.4}\approx\qty{6.5}{m/s} \]
です。雨つぶは、落ち始めてから数秒、十数mのうちにこの速さになり、残りの1000\,m近くを一定の速さで落ちてきます（はかった値：直径2\,mmで約\qty{6.5}{m/s}、5\,mmで約\qty{9}{m/s}。1949年、ガンとキンザーの測定）。

\noindent つぶの大きさで、終端速度はどれくらい変わるのでしょう。はかった値をくらべます。
\begin{bkcenter}
\begin{bktable}{colspec={Q[l,m,wd=27mm] *{5}{Q[c,m,wd=17mm]}}, rowsep=1.5pt, colsep=3pt}
つぶ & 霧 & 霧雨 & 小雨 & 雨 & 大つぶ \\
直径（mm） & 0.02 & 0.2 & 1 & 2 & 5 \\
終端速度（m/s） & 0.012 & 0.7 & 4.0 & 6.5 & 9.1 \\
1000\,mに & 約23時間 & 約24分 & 約4分 & 約2.6分 & 約1.8分 \\
\end{bktable}
\bktabcaption{つぶの大きさと終端速度（霧は⑤で計算する値、ほかはガンとキンザーの測定、1949年）。1000\,m÷終端速度で、落ちる時間を計算した。}
\end{bkcenter}
小さいつぶほど、ずっとゆっくり落ちます。霧のつぶは1000\,m落ちるのに丸1日ほどかかるので、空に浮かんで見えます。

\bkvifhead{重い物は速く落ちる（空気があるとき）}
同じ形・同じ大きさで、質量だけ2倍にすると、
\begin{itemize}[leftmargin=1.2\zw, itemsep=1pt, topsep=2pt, label={\color{\bkTT}●}]
\item 速さに比例する抵抗なら、$v_{\mathrm f}=\dfrac{mg}{b}$ は $m$ に比例するので、2倍。
\item 2乗に比例する抵抗なら、$v_{\mathrm t}=\sqrt{\dfrac{mg}{c}}$ は $\sqrt{m}$ に比例するので、$\sqrt2\approx1.41$ 倍。4倍なら $\sqrt4=2$ 倍。
\end{itemize}
やってみようのフィルターは、重ねても形と大きさがほとんど同じなので、$c$ は同じで、質量だけがふえます。2枚で約1.4倍、4枚で約2倍になったのは、2乗に比例する抵抗だったからです（フィルターの4枚で1.93倍と、2倍より少し小さかったのは、はかったときのずれです）。

空気がなければ（5-13）、重い物も軽い物も同じ速さで落ちます。空気の中で重い物が速く落ちるのは、重い物ほど「抵抗とつり合うまで」速くなれるからです。

\bksubsection{4}{表計算で、落ちる速さを計算する}

\leavevmode\label{bk:6-5:hyou}終端速度に「どのように」近づくかは、6-2のオイラー法で計算できます。スカイダイバー（質量\qty{80}{kg}、手足を広げた姿勢で $c=\qty{0.26}{N.s^2/m^2}$）を考えます。式の両辺を $m$ でわると、
\[ \frac{\dd v}{\dd t}=g-\frac{c}{m}\,v^2=9.8-0.00325\,v^2 \]
（$\dfrac{c}{m}=\dfrac{0.26}{80}=0.00325$）。終端速度は $v_{\mathrm t}=\sqrt{\dfrac{80\times9.8}{0.26}}=\sqrt{3015}\approx\qty{54.9}{m/s}$（時速約198\,km）です。

\noindent $\Delta t=1$ 秒として、6-2と同じ手順でくり返します。
\begin{enumerate}[label={\bkcircled[\bkTF]{\arabic*}}, leftmargin=2.4\zw, itemsep=1pt, topsep=2pt]
\item 今の速さから、今の加速度 $a=9.8-0.00325\,v^2$ を計算する。
\item 次の速さ＝今の速さ＋$a\times\Delta t$。
\end{enumerate}
\begin{bkcenter}
\bksheet{4}{24}{%
1 & $t$（s） & $v$（m/s） & $a$（m/s$^2$） & 空気なし $9.8t$ \\
2 & 0 & 0 & 9.80 & 0 \\
3 & 1 & 9.80 & 9.49 & 9.8 \\
4 & 2 & 19.29 & 8.59 & 19.6 \\
5 & 3 & 27.88 & 7.27 & 29.4 \\
6 & 4 & 35.15 & 5.78 & 39.2 \\
7 & 5 & 40.94 & 4.35 & 49.0 \\
8 & 6 & 45.29 & 3.13 & 58.8 \\
9 & 7 & 48.42 & 2.18 & 68.6 \\
10 & 8 & 50.60 & 1.48 & 78.4 \\
11 & 9 & 52.08 & 0.98 & 88.2 \\
12 & 10 & 53.07 & 0.65 & 98.0 \\}
\bktabcaption{$\Delta t=1$ 秒のオイラー法。C列の式は \texttt{=9.8-0.00325*B2\^{}2}、B列の式は \texttt{=B2+C2*1}。}
\end{bkcenter}
\begin{bkfigure}
\begin{tikzpicture}
\begin{axis}[bkplot, width=108mm, height=52mm, xmin=0, xmax=16.5, ymin=0, ymax=68, xtick={2,4,...,16}, ytick={20,40,54.9}, yticklabels={20,40,54.9}, xlabel={$t$（s）}, ylabel={$v$（m/s）}]
  \addplot[bkthin, dashed, domain=0:16] {54.9};
  \addplot[bkcurve2, line width=1.2pt, domain=0:6.6] {9.8*x};
  \addplot[bkcurve, domain=0:16, samples=80] {54.91*tanh(9.8*x/54.91)};
  \addplot[only marks, mark=*, mark size=1.5pt, bkInk] coordinates {(0,0) (1,9.8) (2,19.29) (3,27.88) (4,35.15) (5,40.94) (6,45.29) (7,48.42) (8,50.60) (9,52.08) (10,53.07)};
  \node[font=\sffamily\footnotesize, text=bkOrange!80!black, anchor=south east, fill=white, inner sep=1pt] at (axis cs:5.3,58.5) {空気がないとき $v=9.8t$};
  \node[font=\sffamily\footnotesize, anchor=south east] at (axis cs:16,56.5) {終端速度 54.9\,m/s};
  \node[font=\sffamily\footnotesize, text=bkBlue, anchor=north west] at (axis cs:10.5,49) {正しい答え（青）と表計算（黒い点）};
\end{axis}
\end{tikzpicture}
\bkfigcaption{スカイダイバーの速さ。はじめは空気がないときとほとんど同じに速くなるが、だんだん速くなり方がにぶくなり、54.9\,m/sに近づく。}
\end{bkfigure}
表計算の値は、正しい値（6-11で式で求める）より少し大きめです。$\Delta t=0.1$ 秒にすると、正しい値とほぼ同じになります。
\begin{bkcenter}
\begin{bktable}{colspec={Q[l,m,wd=40mm] *{5}{Q[c,m,wd=11mm]}}, rowsep=1.5pt}
時間（s） & 1 & 2 & 3 & 5 & 10 \\
オイラー法 $\Delta t=1$ & 9.80 & 19.29 & 27.88 & 40.94 & 53.07 \\
オイラー法 $\Delta t=0.1$ & 9.71 & 18.86 & 26.98 & 39.30 & 52.02 \\
正しい値（6-11） & 9.70 & 18.81 & 26.88 & 39.13 & 51.90 \\
\end{bktable}
\bktabcaption{$\Delta t$ を小さくすると、正しい値に近づく（6-2）。10秒で終端速度の約95\,\%。}\label{bk:6-5:tab-dt}
\end{bkcenter}
はじめの1秒は、空気がないとき（9.8\,m/s）とほとんど同じです。速さが小さいうちは、抵抗が小さいからです。10秒ほどで、ほぼ終端速度に達します。

\begin{seikaku}
2乗に比例する抵抗の式 $m\dfrac{\dd v}{\dd t}=mg-cv^2$ も、式で解くことができます。答えは $v=v_{\mathrm t}\tanh\dfrac{gt}{v_{\mathrm t}}$ で、$\tanh$（ハイパボリック・タンジェント）という関数を使います。6-11で学びます。表\ref{bk:6-5:tab-dt}の「正しい値」は、この式で計算しました。
\end{seikaku}

\bksubsection{5}{速さに比例する抵抗の式を解く}

\leavevmode\label{bk:6-5:toku}\bklevel{3}\enspace 速さに比例する抵抗なら、式は $m\dfrac{\dd v}{\dd t}=mg-bv$ です。これは、6-4の「目標の値へ近づく形」にできます。
\begin{bkhenkei}
m\,\frac{\dd v}{\dd t}=mg-bv & 運動の式（下向きが $+$） \\
\frac{\dd v}{\dd t}=g-\frac{b}{m}\,v & 両辺を $m$ でわる \\
\frac{\dd v}{\dd t}=\frac{b}{m}\Bigl(\frac{mg}{b}-v\Bigr) & 右辺を $\frac{b}{m}$ でくくる（$g=\frac{b}{m}\times\frac{mg}{b}$） \\
\frac{\dd v}{\dd t}=\frac{v_{\mathrm f}-v}{\tau} & $v_{\mathrm f}=\frac{mg}{b}$、$\tau=\frac{m}{b}$ とおく \\
\end{bkhenkei}
最後の行は、6-4④の形そのものです（目標の値が $v_{\mathrm f}$）。はじめ止まっていた（$v_0=0$）とすると、6-4の答えの形 $y=y_\infty+(y_0-y_\infty)e^{-t/\tau}$ に、$y_\infty=v_{\mathrm f}$、$y_0=0$ を入れて、
\bkvifframe{$v=v_{\mathrm f}\,(1-e^{-t/\tau})$\qquad $v_{\mathrm f}=\dfrac{mg}{b}$、\quad $\tau=\dfrac{m}{b}$}
これは、物理シリーズ『力学』の力-2で予告した答えです。6-4で学んだように、$\tau$ で終端速度の63\,\%、$3\tau$ で95\,\%、$5\tau$ で99.3\,\%に達します。

\bkvifhead{変数分離でも解ける}
終端速度との差 $u=v_{\mathrm f}-v$ を主役にします（$v_{\mathrm f}$ は定数なので $\frac{\dd u}{\dd t}=-\frac{\dd v}{\dd t}$）。
\begin{bkhenkei}
\frac{\dd u}{\dd t}=-\frac{u}{\tau} & 上の式の両辺に $-1$ をかけ、$v_{\mathrm f}-v=u$ とおく \\
\int\frac{\dd u}{u}=\int\Bigl(-\frac{1}{\tau}\Bigr)\dd t & 変数分離（6-4③）：両辺を $u$ でわって積分 \\
\ln u=-\frac{t}{\tau}+C & 左辺は $\ln u$（$u>0$） \\
u=Ae^{-t/\tau} & $e$ の○乗にもどし、$e^C=A$ とおく \\
A=v_{\mathrm f} & $t=0$ で $v=0$、つまり $u=v_{\mathrm f}$ \\
v=v_{\mathrm f}-v_{\mathrm f}e^{-t/\tau} & $v=v_{\mathrm f}-u$ にもどす \\
\end{bkhenkei}
同じ答え $v=v_{\mathrm f}(1-e^{-t/\tau})$ になりました。

\bkvifhead{はじめは $v\approx gt$ になっている？}
$t$ が $\tau$ よりずっと小さいとき、5-8の近似 $e^{-x}\approx1-x$ を使うと、
\[ v=v_{\mathrm f}\bigl(1-e^{-t/\tau}\bigr)\approx v_{\mathrm f}\Bigl(1-\Bigl(1-\frac{t}{\tau}\Bigr)\Bigr)=v_{\mathrm f}\times\frac{t}{\tau}=\frac{mg}{b}\times\frac{b}{m}\,t=gt \]
落ち始めは、空気がないとき（$v=gt$）と同じです。抵抗がまだ小さいからで、図の表計算とも合っています。

\bkvifhead{$\tau$ の意味と単位}
$\tau=\dfrac{m}{b}$ の単位は、$\dfrac{\text{kg}}{\text{N·s/m}}=\dfrac{\text{kg·m}}{\text{N·s}}$ です。$\text{N}=\text{kg·m/s}^2$（$F=ma$ から）なので、
\[ \frac{\text{kg·m}}{\text{kg·m/s}^2\times\text{s}}=\frac{1}{1/\text{s}}=\text{s} \]
たしかに時間です。また、$\tau=\dfrac{m}{b}=\dfrac{mg}{b}\div g=\dfrac{v_{\mathrm f}}{g}$ なので、$\tau$ は「加速度 $g$ のまま速くなったら、終端速度に着く時間」です（6-4の「はじめの接線が目標に届く時間」）。

\bkvifhead{質量や抵抗の係数を2倍にすると}
$v_{\mathrm f}=\dfrac{mg}{b}$ と $\tau=\dfrac{m}{b}$ の式から、答えの曲線がどう変わるかを読みます。
\begin{itemize}[leftmargin=1.2\zw, itemsep=1pt, topsep=2pt, label={\color{\bkTT}●}]
\item 質量 $m$ を2倍：$v_{\mathrm f}$ も $\tau$ も2倍。速く落ちるが、そこに近づくのに2倍の時間がかかる。
\item 抵抗の係数 $b$ を2倍：$v_{\mathrm f}$ も $\tau$ も $\frac{1}{2}$ 倍。おそい速さに、すばやく落ち着く。
\end{itemize}
\begin{bkfigure}
\begin{tikzpicture}
\begin{axis}[bkplot, width=104mm, height=46mm, xmin=0, xmax=8.6, ymin=0, ymax=2.35, xtick={1,2,4,6,8}, ytick={0.5,1,2}, xlabel={$t$}, ylabel={$v$}]
  \addplot[bkcurve, domain=0:8.3, samples=60] {1-exp(-x)};
  \addplot[bkcurve2, domain=0:8.3, samples=60] {2*(1-exp(-x/2))};
  \addplot[bkcurve3, domain=0:8.3, samples=60] {0.5*(1-exp(-2*x))};
  \addplot[bkthin, dashed, domain=0:2.1] {x};
  \node[font=\sffamily\footnotesize, text=bkOrange!80!black, anchor=south east, fill=white, inner sep=1pt] at (axis cs:8.3,2.0) {$m$ を2倍};
  \node[font=\sffamily\footnotesize, text=bkBlue, anchor=south east, fill=white, inner sep=1pt] at (axis cs:8.3,1.05) {もと};
  \node[font=\sffamily\footnotesize, text=bkGreen!60!black, anchor=north east, fill=white, inner sep=1pt] at (axis cs:8.3,0.45) {$b$ を2倍};
  \node[font=\sffamily\footnotesize, text=bkGray, anchor=west, fill=white, inner sep=1pt] at (axis cs:2.15,2.1) {$v=gt$};
\end{axis}
\end{tikzpicture}
\bkfigcaption{速さに比例する抵抗で落ちる物の速さ（横軸は「もとの $\tau$ の何倍か」、たて軸は「もとの $v_{\mathrm f}$ の何倍か」）。どの曲線も、落ち始めは灰色の点線 $v=gt$（空気がないとき）にそって速くなる。}
\end{bkfigure}
どの曲線も、はじめは同じ傾き $g$ で始まります。落ち始めは抵抗が小さく、どれも空気がないときと同じに速くなるからです（「はじめは $v\approx gt$」）。

\begin{motto}[3]
\textbf{霧のつぶの $\frac{b}{m}$ を計算する（ストークスの法則）}\par
小さな球が、ねばりけのある空気や液体の中をゆっくり動くとき、抵抗の係数は
\[ b=6\pi\eta r \]
になることを、イギリスのストークスが1851年に示しました（ストークスの法則）。\bkidx{すとーくすのほうそく}{ストークスの法則}$r$ は球の半径、$\eta$（イータ）は空気のねばりけ（粘度）で、約 $1.8\times10^{-5}$\,Pa·s です。水の球の質量は、密度 $\rho=\qty{1000}{kg/m^3}$ に球の体積 $\frac{4}{3}\pi r^3$ をかけたものです。
\begin{bkhenkei}
\frac{b}{m}=\frac{6\pi\eta r}{\rho\times\frac{4}{3}\pi r^3} & $b$ と $m$ を入れた \\
\frac{b}{m}=\frac{9\eta}{2\rho r^2} & $\pi$ と $r$ を約分（$6\div\frac{4}{3}=\frac{9}{2}$） \\
\frac{b}{m}=\frac{9\times1.8\times10^{-5}}{2\times1000\times(10^{-5})^2}=810 & $r=\qty{10}{\micro m}=10^{-5}$\,m を入れた \\
\end{bkhenkei}
$r^2$ が分母にあるので、半径が半分なら $\frac{b}{m}$ は4倍、終端速度は $\frac{1}{4}$ です。
\end{motto}

\bkvifhead{霧のつぶが浮かんで見えるわけ}
半径0.01\,mm（\qty{10}{\micro m}）の霧のつぶは、空気のねばりけによる抵抗を受けて落ちるので、速さに比例する抵抗になります。このつぶでは $\dfrac{b}{m}=810$（1/s）です（つぶの大きさと空気のねばりけから計算した値。すぐ前のもっと知りたい）。
\[ v_{\mathrm f}=\frac{mg}{b}=\frac{g}{b/m}=\frac{9.8}{810}\approx\qty{0.012}{m/s},\qquad \tau=\frac{m}{b}=\frac{1}{810}\approx\qty{0.0012}{s} \]
1000分の1秒ほどで終端速度になり、1秒に1.2\,cmしか落ちません。1\,m落ちるのに80秒以上かかります。わずかな上向きの風でも落ちなくなるので、霧や雲は空に浮かんで見えるのです。




\bksubsection{6}{エレベーターの中の体重計}

\leavevmode\label{bk:6-5:elevator}5-5で「くわしくは第6巻で」と予告した、エレベーターの中の体重計を、運動の式で調べます。ここでは\textbf{上向きを $+$} にします（エレベーターが上へ動くので）。
\par\vspace{3pt}\noindent\begin{minipage}[c]{0.6\linewidth}
体重計に乗った人（質量 $m$）にはたらく力は、
\begin{itemize}[leftmargin=1.2\zw, itemsep=1pt, topsep=2pt, label={\color{\bkTT}●}]
\item 重力 $mg$（下向き）
\item 体重計が足をおし上げる力 $N$（上向き）
\end{itemize}
の2つです。人はエレベーターといっしょに動くので、人の加速度はエレベーターの加速度 $a$ と同じです。上向きの力の合計は $N-mg$ なので、運動の法則から
\[ m\,\frac{\dd v}{\dd t}=N-mg \]
\end{minipage}\hfill
\begin{minipage}[c]{0.36\linewidth}\centering
\begin{tikzpicture}[x=1mm, y=1mm]
  \draw[bkInk!70, line width=0.9pt] (0,0) rectangle (34,46);
  \draw[bkInk!50, line width=0.6pt] (17,46) -- (17,54);
  \filldraw[fill=bkGray!30, draw=bkInk!60, line width=0.6pt] (9,0.6) rectangle (25,3.4);
  \bkperson{(17,3.4)}{30}{bkBlue}
  \draw[bkRed, line width=1.4pt, -{Stealth[length=4pt]}] (29,22) -- (29,12);
  \node[font=\sffamily\footnotesize, text=bkRed, anchor=west] at (30,14) {$mg$};
  \draw[bkBlue, line width=1.4pt, -{Stealth[length=4pt]}] (5,8) -- (5,22);
  \node[font=\sffamily\footnotesize, text=bkBlue, anchor=east] at (3.8,18) {$N$};
  \draw[bkInk!60, line width=0.8pt, -{Stealth[length=4pt]}] (40,30) -- (40,42);
  \node[font=\sffamily\footnotesize, anchor=north] at (40,29) {$a$};
\end{tikzpicture}
\end{minipage}
\par\vspace{3pt}
\noindent 体重計は、足をおす力 $N$ をはかり、それを $g$ でわって「kg」で表示します。止まっているとき（$a=0$）は $N=mg$ なので、ふだんの体重 $m$ を表示します。
\begin{bkhenkei}
m\,a=N-mg & 加速度 $a=\frac{\dd v}{\dd t}$ \\
N=m\,(g+a) & $mg$ を左へ移す \\
\frac{N}{g}=m\times\frac{g+a}{g} & 体重計の目もり（両辺を $g$ でわる） \\
\end{bkhenkei}
60\,kgの人の目もりは、上へ動き出すとき（$a=\qty{1}{m/s^2}$）と止まるとき（$a=\qty{-1}{m/s^2}$）で、
\[ 60\times\frac{9.8+1}{9.8}=66.1\ \text{kg},\qquad 60\times\frac{9.8-1}{9.8}=53.9\ \text{kg} \]
です。体重計が表示するのは、速さではなく加速度で決まる量です。5-5で「体が感じていたのは加速度だった」と書いたのは、このことです。目もりを時刻ごとにグラフにすると、図\ref{bk:6-5:fig-elev}のようになります。
\begin{bkfigure}
\begin{tikzpicture}
\begin{axis}[bkplot, name=ev, width=100mm, height=25mm, xmin=0, xmax=11.2, ymin=0, ymax=2.7, xtick={2,4,6,8,10}, ytick={1,2}, xlabel={$t$（s）}, ylabel={$v$（m/s）}]
  \addplot[bkcurve, line width=1.4pt] coordinates {(0,0) (2,2) (8,2) (10,0)};
\end{axis}
% Ver2 指摘6：たてを高くし、「重い」「軽い」を線・軸・軸の名前からはなした
\begin{axis}[bkplot, at={(0mm,-46mm)}, width=100mm, height=38mm, xmin=0, xmax=11.2, ymin=44, ymax=75, xtick={2,4,6,8,10}, ytick={53.9,60,66.1}, xlabel={$t$（s）}, ylabel={目もり（kg）}]
  \addplot[bkcurve2, line width=1.4pt, const plot] coordinates {(0,66.1) (2,60) (8,53.9) (10,53.9)};
  \draw[bkGray, dashed, line width=0.5pt] (axis cs:0,60) -- (axis cs:10.6,60);
  \node[font=\sffamily\footnotesize] at (axis cs:1.2,70.3) {重い};
  \node[font=\sffamily\footnotesize] at (axis cs:9,49.4) {軽い};
\end{axis}
\end{tikzpicture}
\bkfigcaption{上：上へ動くエレベーターの速さ（5-5と同じ）。下：60\,kgの人の体重計の目もり（0〜2秒は66.1\,kg、2〜8秒は60\,kg、8〜10秒は53.9\,kg）。}\label{bk:6-5:fig-elev}
\end{bkfigure}

逆に、体重計の目もりを時間ごとに記録しておけば、$a=\dfrac{N}{m}-g$ から加速度がわかり、それを積分すれば（6-3）速さが、もう1回積分すれば位置がわかります。スマートフォンの中の加速度センサーも、このしくみで動きをはかっています。

\bksubsection{7}{電気の終端速度：RL 回路}

\leavevmode\label{bk:6-5:rl}6-4⑦で、電池 $E$、抵抗 $R$、コイル $L$ を直列にした回路の式 $L\dfrac{\dd i}{\dd t}=E-Ri$ を作りました。速さに比例する抵抗を受けて落ちる物の式と、ならべてみます。
\begin{bkfigure}
\begin{tikzpicture}[x=1mm, y=1mm]
  % 左：落ちる物
  % Ver2：図を 5 mm 低くした（ページの下の空き・偶数ページのため。中身は同じ）
  \bkvifdrop{(0,26)}{2.4}
  \draw[bkRed, line width=1.4pt, -{Stealth[length=4pt]}] (0,20) -- (0,10);
  \node[font=\sffamily\footnotesize, text=bkRed, anchor=west] at (1.5,12) {$mg$};
  \draw[bkBlue, line width=1.4pt, -{Stealth[length=4pt]}] (0,33) -- (0,40);
  \node[font=\sffamily\footnotesize, text=bkBlue, anchor=west] at (1.5,38) {$bv$};
  % 右：RL 回路
  \begin{scope}[shift={(54,10)}]
    \draw[bkInk, line width=0.7pt] (0,0) rectangle (30,26);
    \bkvifE{(0,13)}{-90}
    \bkvifR{(15,26)}{0}
    \bkvifL{(30,13)}{-90}
    \bkvifS{(15,0)}{0}
    \node[font=\sffamily\footnotesize, anchor=east] at (-5.5,13) {$E$};
    \node[font=\sffamily\footnotesize, anchor=south] at (15,29.5) {$R$};
    \node[font=\sffamily\footnotesize, anchor=west] at (35,13) {$L$};
    \draw[bk arrow, bkRed, line width=1pt] (2,30) -- (9,30);
    \node[font=\sffamily\footnotesize, text=bkRed, anchor=south] at (5.5,31) {$i$};
  \end{scope}
  \node[bk node, draw=bkBlue, fill=bkBlue!6, minimum width=50mm, minimum height=11mm, inner ysep=1.5mm, font=\small] (a) at (0,0) {$m\,\dfrac{\dd v}{\dd t}=mg-b\,v$};
  \node[bk node, draw=bkOrange, fill=bkOrange!8, minimum width=50mm, minimum height=11mm, inner ysep=1.5mm, font=\small] (b) at (69,0) {$L\,\dfrac{\dd i}{\dd t}=E-R\,i$};
  \draw[bk link, {Stealth[length=4pt]}-{Stealth[length=4pt]}, bkGray] (a) -- (b);
\end{tikzpicture}
\bkfigcaption{左：速さに比例する抵抗を受けて落ちる物。右：RL 回路。2つの式は、文字を入れかえると同じ形になる。}
\end{bkfigure}
\begin{bkcenter}
\begin{bktable}{colspec={Q[l,m,wd=30mm] Q[c,m,wd=26mm] Q[c,m,wd=34mm] X[l,m]}, width=\linewidth, rowsep=1.5pt}
役わり & 落ちる物 & RL 回路 & 意味 \\
変わる量 & 速さ $v$ & 電流 $i$ & 流れの速さ \\
動きにくさ（慣性） & 質量 $m$ & インダクタンス $L$ & 急には変えられない \\
おす力 & 重力 $mg$ & 電池の電圧 $E$ & 一定でおし続ける \\
じゃまするもの & 抵抗の係数 $b$ & 抵抗 $R$ & 流れに比例してじゃまする \\
最後の値 & $v_{\mathrm f}=\dfrac{mg}{b}$ & $\dfrac{E}{R}$ & おす力＝じゃま \\
時定数 & $\dfrac{m}{b}$ & $\dfrac{L}{R}$ & 慣性÷じゃま \\
\end{bktable}
\bktabcaption{落ちる物と RL 回路の対応。}
\end{bkcenter}
RL 回路の最後の電流 $\dfrac{E}{R}$ は、「電流の終端速度」です。スイッチを入れた直後は、コイルが電流の急な変化をきらうので、電流は0から始まります。落ちる物が、はじめ止まっていて、急には速くなれないのと同じです。電流がふえると抵抗の電圧 $Ri$ がふえ、電池の電圧 $E$ と同じになったところで、電流は変わらなくなります。

このように、ちがう物でも式の形が同じなら、同じ答えになります。この対応は、6-15で、ばね・電気・熱をまとめるときにもう一度使います。



\begin{column}{雨つぶの形は「しずく形」ではない}
絵本などでは、雨つぶは上がとがった「しずく形」にかかれます。でも、空気の中を落ちている雨つぶは、小さいものは丸く、直径2\,mmをこえると、下からの空気におされて、下が平らな「まんじゅう形」になります。もっと大きくなると、下がへこんでパラシュートのような形になり、やがて分かれて小さなつぶになります。だから、とても大きな雨つぶはなく、終端速度も毎秒9\,mほどでそれ以上大きくなりません。
\end{column}

\begin{senmonsho}[大学の教科書ふうの書き方]
\begin{bkhonbun}
速度に比例する抵抗を受けて鉛直に落下する質点の運動方程式は $m\dot v=mg-bv$ であり、初速度0の解は
\[ v(t)=v_{\mathrm f}\bigl(1-e^{-t/\tau}\bigr),\qquad v_{\mathrm f}=mg/b,\quad \tau=m/b \]
である。
\end{bkhonbun}
\bkyaku{「鉛直」は真下の向き。「質点」は大きさを考えない物。$\dot v$ は $\frac{\dd v}{\dd t}$（5-5）。「初速度」は、はじめの速度。}
\end{senmonsho}

\begin{taisetsu}
\def\tsl#1{\makebox[6.6\zw][l]{\sffamily\bfseries\small\color{\bkTT}#1}}
\tsl{落ちる物の式}$m\dfrac{\dd v}{\dd t}=mg-f$\quad（下向きが $+$、抵抗 $f$ は動きと反対）\\[4pt]
\tsl{終端速度}$\dfrac{\dd v}{\dd t}=0$ とおく：$v_{\mathrm f}=\dfrac{mg}{b}$（$f=bv$）、$v_{\mathrm t}=\sqrt{\dfrac{mg}{c}}$（$f=cv^2$）\\[4pt]
\tsl{速さの式}$f=bv$ なら $v=v_{\mathrm f}(1-e^{-t/\tau})$、$\tau=\dfrac{m}{b}$\\[4pt]
\tsl{体重計}$N=m(g+a)$\quad（上向きが $+$）
\end{taisetsu}

\begin{kotoba}
\small
\bkword{抵抗の係数}{空気などの抵抗の大きさを決める数。$f=bv$ の $b$〔N·s/m〕、$f=cv^2$ の $c$〔N·s$^2$/m$^2$〕。}
\bkword{終端速度}{抵抗が重力とつり合って、それ以上速くならなくなった速さ。}
\bkword{運動の式（運動方程式）}{質量×加速度＝力の合計。$\frac{\dd v}{\dd t}$ が入る微分方程式。}
\end{kotoba}

\bkkigocards{\bkkigocard{$v_{\mathrm f}$、$v_{\mathrm t}$}{ブイ・エフ、ブイ・ティー}{終端速度。$v_{\mathrm f}$ は速さに比例する抵抗、$v_{\mathrm t}$ は速さの2乗に比例する抵抗のとき}{terminal velocity}{6-5}}
{\bkkigocard{$b$、$c$}{ビー、シー}{抵抗の係数（$f=bv$、$f=cv^2$）。6-7のダンパーの $c$ とは別の量}{drag coefficient}{6-5}}

%% ---------------- 例題 ----------------
\bksection{reidai}{例題でたしかめよう}

\begin{reidai}[速さに比例する抵抗の終端速度]
半径\qty{10}{\micro m}の霧のつぶは、$\dfrac{\dd v}{\dd t}=9.8-810\,v$（$v$ はm/s、$t$ は秒）にしたがって落ちます。終端速度、時定数、終端速度で1\,m落ちる時間を求めましょう。\par
\kangaekata 終端速度は $\dfrac{\dd v}{\dd t}=0$ とおく（③）。$\tau=\dfrac{m}{b}$ は、$v$ の係数810（$=\dfrac{b}{m}$）の逆数。\par
\shiki $0=9.8-810\,v_{\mathrm f}$ より $v_{\mathrm f}=\dfrac{9.8}{810}=0.0121$（m/s）\bkvifq
$\tau=\dfrac{1}{810}=0.00123$（s）\bkvifq
1\,m落ちる時間：$\dfrac{\qty{1}{m}}{\qty{0.0121}{m/s}}=\qty{82.6}{s}$\par
\kotae 約\qty{0.012}{m/s}（1.2\,cm/s）、約0.0012秒、約83秒
\end{reidai}

\begin{jibun}
半径\qty{5}{\micro m}のもっと小さな霧のつぶは、$\dfrac{\dd v}{\dd t}=9.8-3240\,v$ にしたがって落ちます。終端速度と時定数を求めましょう。また、終端速度で1\,m落ちるのに何秒かかりますか。
\bkjans{$v_{\mathrm f}=\dfrac{9.8}{3240}\approx\qty{0.0030}{m/s}$。$\tau=\dfrac{1}{3240}\approx\qty{0.00031}{s}$。1\,mに約330秒（$1\div0.003025\approx331$）。}
\end{jibun}

\begin{reidai}[2乗に比例する抵抗の終端速度]
質量\qty{80}{kg}のスカイダイバーが、手足を広げて落ちています。抵抗は $f=cv^2$、$c=\qty{0.26}{N.s^2/m^2}$ です。終端速度は何m/sですか。時速何kmですか。\par
\kangaekata ③の $v_{\mathrm t}=\sqrt{\dfrac{mg}{c}}$。時速にするには \qty{3600}{s/h}、\qty{1000}{m/km} を使う。\par
\shiki $v_{\mathrm t}=\sqrt{\dfrac{80\times9.8}{0.26}}=\sqrt{\dfrac{784}{0.26}}=\sqrt{3015.4}=54.9$\bkvifq
$\qty{54.9}{m/s}\times\qty{3600}{s/h}=\qty{197640}{m/h}=\qty{197.6}{km/h}$\par
\kotae 約\qty{54.9}{m/s}、時速約198\,km
\end{reidai}

\begin{jibun}
例題2のスカイダイバーがパラシュートを開くと、$c=\qty{30}{N.s^2/m^2}$ になりました（質量はそのまま\qty{80}{kg}）。終端速度は何m/sですか。
\bkjans{$v_{\mathrm t}=\sqrt{\dfrac{80\times9.8}{30}}=\sqrt{26.1}\approx\qty{5.1}{m/s}$。}
\end{jibun}

\begin{reidai}[質量を変えると]
形と大きさが同じで、2乗に比例する抵抗を受けて落ちる物があります。質量を4倍にすると、終端速度は何倍になりますか。また、質量を2倍にすると何倍ですか（$\sqrt2=1.41$）。\par
\kangaekata $v_{\mathrm t}=\sqrt{\dfrac{mg}{c}}$ で、$g$ と $c$（形と大きさで決まる）は変わらない。$v_{\mathrm t}$ は $\sqrt{m}$ に比例する。\par
\shiki 4倍：$\sqrt{\dfrac{4m\times g}{c}}=\sqrt4\times\sqrt{\dfrac{mg}{c}}=2\times v_{\mathrm t}$\bkvifq
2倍：$\sqrt{\dfrac{2m\times g}{c}}=\sqrt2\times\sqrt{\dfrac{mg}{c}}=1.41\times v_{\mathrm t}$\par
\kotae 4倍で2倍、2倍で約1.41倍
\end{reidai}
\noindent\begin{minipage}[c]{0.5\linewidth}
\begin{tikzpicture}
\begin{axis}[bkplot, width=60mm, height=42mm, xmin=0, xmax=4.9, ymin=0, ymax=2.6, xtick={1,2,3,4}, ytick={1,2}, xlabel={枚数}, ylabel={速さ（m/s）}]
  \addplot[bkcurve, domain=0:4.6, samples=60] {1.06*sqrt(x)};
  \addplot[only marks, mark=*, mark size=1.8pt, bkOrange] coordinates {(1,1.06) (2,1.47) (4,2.05)};
\end{axis}
\end{tikzpicture}
\end{minipage}\hfill
\begin{minipage}[c]{0.46\linewidth}\raggedright
やってみようのフィルターの結果（オレンジの点）を、$1.06\sqrt{\text{枚数}}$ の曲線（青）とくらべます。点は曲線の近くに並びます。終端速度は質量の $\sqrt{\ }$ に比例するので、質量をふやしても、速さはゆっくりとしかふえません。
\end{minipage}
\par\vspace{4pt}

\begin{jibun}
例題3と同じように、形と大きさが同じで、2乗に比例する抵抗を受けて落ちる物があります。質量を9倍にすると、終端速度は何倍になりますか。また、質量を16倍にすると何倍ですか。
\bkjans{$\sqrt9=3$ 倍、$\sqrt{16}=4$ 倍。}
\end{jibun}

\begin{reidai}[表計算の1行を手で計算する]
スカイダイバーの速さが $\dfrac{\dd v}{\dd t}=9.8-0.00325\,v^2$ にしたがいます。はじめ $v=0$ として、$\Delta t=1$ 秒のオイラー法で、1秒後と2秒後の速さを求めましょう。\par
\kangaekata ④の手順。今の速さから加速度を計算し、次の速さ＝今の速さ＋加速度×$\Delta t$。\par
\shiki 0秒：加速度 $9.8-0.00325\times0^2=9.8$\bkvifq
1秒後の速さ：$0+9.8\times1=9.8$\bkvifq
1秒：加速度 $9.8-0.00325\times9.8^2=9.8-0.00325\times96.04$\bkvifq
\hspace*{3\zw}$=9.8-0.312=9.488$\bkvifq
2秒後の速さ：$9.8+9.488\times1=19.288$\par
\kotae 1秒後\qty{9.8}{m/s}、2秒後約\qty{19.3}{m/s}
\end{reidai}
\noindent\begin{minipage}[c]{0.5\linewidth}
\begin{tikzpicture}
\begin{axis}[bkplot, width=60mm, height=42mm, xmin=0, xmax=8.8, ymin=0, ymax=56, xtick={2,4,6,8}, ytick={20,40}, xlabel={$t$（s）}, ylabel={$v$（m/s）}]
  \addplot[bkcurve, domain=0:8.4, samples=60] {54.91*tanh(9.8*x/54.91)};
  \addplot[bkcurve2, line width=1pt, mark=*, mark size=1.2pt] coordinates {(0,0) (1,9.8) (2,19.29) (3,27.88) (4,35.15) (5,40.94) (6,45.29) (7,48.42) (8,50.60)};
\end{axis}
\end{tikzpicture}
\end{minipage}\hfill
\begin{minipage}[c]{0.46\linewidth}\raggedright
オイラー法（オレンジ）は、1秒ごとに「その時刻の加速度」の傾きでまっすぐ進む折れ線です。1秒のあいだにも加速度は小さくなっていくので、正しい答え（青）より少し上にずれます。$\Delta t$ を小さくすると、ずれも小さくなります（6-2）。
\end{minipage}
\par\vspace{4pt}

\begin{jibun}
バドミントンのシャトルの速さが $\dfrac{\dd v}{\dd t}=9.8-0.2\,v^2$ にしたがいます。はじめ $v=0$ として、$\Delta t=0.2$ 秒のオイラー法で、0.2秒後と0.4秒後の速さを求めましょう。また、終端速度を求めましょう。
\bkjans{0.2秒後 $0+9.8\times0.2=\qty{1.96}{m/s}$。加速度 $9.8-0.2\times1.96^2\approx9.032$、0.4秒後 $1.96+9.032\times0.2\approx\qty{3.77}{m/s}$。終端速度 $\sqrt{\dfrac{9.8}{0.2}}=\sqrt{49}=\qty{7}{m/s}$。}
\end{jibun}

\begin{reidai}[速さの式を求める]
例題1の霧のつぶ（$\dfrac{\dd v}{\dd t}=9.8-810\,v$）が、止まった状態から落ち始めます。$v$ を $t$ の式で表し、$3\tau$ たったときの速さを求めましょう（$e^{-3}=0.050$）。\par
\kangaekata ⑤の答え $v=v_{\mathrm f}(1-e^{-t/\tau})$ に、例題1の $v_{\mathrm f}$ と $\tau$ を入れる。\par
\shiki $v_{\mathrm f}=0.0121$、$\tau=\dfrac{1}{810}$ なので $\dfrac{t}{\tau}=810\,t$\bkvifq
$v=0.0121\,(1-e^{-810t})$\bkvifq
$3\tau$ では $e^{-3}$ なので、$v=0.0121\times(1-0.050)$\bkvifq
\hspace*{3\zw}$=0.0121\times0.95=0.0115$\par
\kotae $v=0.0121(1-e^{-810t})$、$3\tau$（約0.0037秒）で約\qty{0.0115}{m/s}（終端速度の95\,\%）
\end{reidai}
\noindent\begin{minipage}[c]{0.5\linewidth}
\begin{tikzpicture}
\begin{axis}[bkplot, width=60mm, height=42mm, xmin=0, xmax=6.6, ymin=0, ymax=14.5, xtick={1.23,3.7,6}, xticklabels={$\tau$,$3\tau$,6}, ytick={11.5,12.1}, yticklabels={11.5,\phantom{1}}, xlabel={$t$（ms）}, ylabel={$v$（mm/s）}]
  \addplot[bkthin, dashed, domain=0:6.3] {12.1};
  \addplot[bkcurve, domain=0:6.3, samples=60] {12.1*(1-exp(-x/1.2346))};
  \draw[bkGray, dashed, line width=0.5pt] (axis cs:3.7,0) -- (axis cs:3.7,11.5);
  \node[font=\sffamily\footnotesize, anchor=south east] at (axis cs:6.3,12.5) {12.1（終端速度）};
\end{axis}
\end{tikzpicture}
\end{minipage}\hfill
\begin{minipage}[c]{0.46\linewidth}\raggedright
横軸をミリ秒（ms、1000分の1秒）、たて軸をmm/sにしてかきました。単位をかえると、6-4の「目標に近づく曲線」と同じ形です。3.7\,msで95\,\%。人の目には、はじめから一定の速さで落ちているように見えます。
\end{minipage}
\par\vspace{4pt}

\begin{jibun}
じぶんでとこう1の霧のつぶ（$\dfrac{\dd v}{\dd t}=9.8-3240\,v$）が、止まった状態から落ち始めます。$v$ を $t$ の式で表し、$3\tau$ たったときの速さを求めましょう（$e^{-3}=0.050$）。
\bkjans{$v=0.00302\,(1-e^{-3240t})$。$3\tau$（約0.00093秒）で $0.00302\times0.95\approx\qty{0.0029}{m/s}$。}
\end{jibun}

\begin{reidai}[エレベーターの体重計]
体重60\,kgの人が、エレベーターの中で体重計に乗っています。エレベーターが（1）上向きに $\qty{1}{m/s^2}$ で加速するとき、（2）上向きに動いていて、$\qty{1}{m/s^2}$ の大きさで減速する（加速度 $\qty{-1}{m/s^2}$）とき、体重計の目もりは何kgですか（$g=9.8$）。\par
\kangaekata ⑥の $\dfrac{N}{g}=m\times\dfrac{g+a}{g}$。上向きが $+$。減速するときは、加速度が下向き（$-$）。\par
\shiki（1）$60\times\dfrac{9.8+1}{9.8}=60\times\dfrac{10.8}{9.8}=66.1$\bkvifq
（2）$60\times\dfrac{9.8+(-1)}{9.8}=60\times\dfrac{8.8}{9.8}=53.9$\par
\kotae（1）約66.1\,kg、（2）約53.9\,kg
\end{reidai}
\noindent\begin{minipage}[c]{0.5\linewidth}
\begin{tikzpicture}
% Ver2 指摘2：たて軸を左はしに置き、目もりの数（53.9・66.1）を点と線からはなした
\begin{axis}[bkplot, axis y line=left, width=60mm, height=50mm, xmin=-10.8, xmax=4.4, ymin=0, ymax=92, xtick={-9.8,-1,1}, xticklabels={$-g$,$-1$,1}, ytick={53.9,66.1}, xlabel={$a$}, ylabel={目もり（kg）}]
  \addplot[bkcurve2, domain=-9.8:4] {60*(9.8+x)/9.8};
  \addplot[only marks, mark=*, mark size=1.6pt, bkInk] coordinates {(1,66.1) (-1,53.9) (-9.8,0)};
  \draw[bkGray, dashed, line width=0.5pt] (axis cs:0,60) -- (axis cs:4,60);
\end{axis}
\end{tikzpicture}
\end{minipage}\hfill
\begin{minipage}[c]{0.46\linewidth}\raggedright
60\,kgの人の目もりと、加速度 $a$ の関係です。$a$ が1ふえるごとに約6.1\,kgふえる直線です。もしケーブルが切れて $a=-g$（自由落下）になると、目もりは0です。体重計の上で「重さがなくなった」ように感じます（無重量）。
\end{minipage}
\par\vspace{4pt}

\begin{jibun}
体重50\,kgの人が、エレベーターの中で体重計に乗っています。加速度が（1）上向きに $\qty{1.5}{m/s^2}$、（2）下向きに $\qty{1.5}{m/s^2}$ のとき、目もりは何kgですか（$g=9.8$）。
\bkjans{（1）\nobreak $50\times\dfrac{11.3}{9.8}\approx57.7$\,kg。（2）\nobreak $50\times\dfrac{8.3}{9.8}\approx42.3$\,kg。}
\end{jibun}

\begin{reidai}[体重計の目もりから加速度を求める]
体重60\,kgの人が、エレベーターの中で体重計に乗ると、目もりが63.0\,kgになりました。エレベーターの加速度の向きと大きさを求めましょう（$g=9.8$）。\par
\kangaekata ⑥の $\dfrac{N}{g}=m\times\dfrac{g+a}{g}$ を、$a$ について解く（逆向きの計算）。上向きが $+$。\par
\shiki $63.0=60\times\dfrac{9.8+a}{9.8}$\bkvifq
両辺を60でわる：$\dfrac{63.0}{60}=1.05=\dfrac{9.8+a}{9.8}$\bkvifq
両辺に9.8をかける：$9.8\times1.05=10.29=9.8+a$\bkvifq
9.8を右辺から移す：$a=10.29-9.8=0.49$\par
\kotae 上向きに約\qty{0.49}{m/s^2}（上へ動き出すとき、または下へ動いていて止まるとき）
\end{reidai}

\begin{jibun}
体重50\,kgの人が、エレベーターの中で体重計に乗ると、目もりが46.0\,kgになりました。エレベーターの加速度の向きと大きさを求めましょう（$g=9.8$）。
\bkjans{$\dfrac{46}{50}=0.92=\dfrac{9.8+a}{9.8}$ より $a=9.8\times0.92-9.8\approx-0.78$。下向きに約\qty{0.78}{m/s^2}。}
\end{jibun}

\begin{reidai}[RL 回路の「終端速度」]
6-4のリレーのコイル（$E=\qty{12}{V}$、$R=\qty{400}{\ohm}$、$L=\qty{0.2}{H}$）で、（1）最後の電流（電流の終端速度）、（2）時定数、（3）スイッチを入れて $3\tau$ たったときの電流を求めましょう（$e^{-3}=0.050$）。\par
\kangaekata ⑦の対応で、$v_{\mathrm f}=\dfrac{mg}{b}$ が $\dfrac{E}{R}$、$\tau=\dfrac{m}{b}$ が $\dfrac{L}{R}$。答えの形も $i=\dfrac{E}{R}(1-e^{-t/\tau})$。\par
\shiki（1）$\dfrac{\qty{12}{V}}{\qty{400}{\ohm}}=\qty{0.03}{A}=\qty{30}{mA}$\bkvifq
（2）$\dfrac{\qty{0.2}{H}}{\qty{400}{\ohm}}=\qty{0.0005}{s}=\qty{0.5}{ms}$\bkvifq
（3）$30\times(1-e^{-3})=30\times0.95=28.5$\par
\kotae（1）\qty{30}{mA}、（2）\qty{0.5}{ms}、（3）約\qty{28.5}{mA}（1.5\,ms後）
\end{reidai}
\noindent\begin{minipage}[c]{0.5\linewidth}
\begin{tikzpicture}
\begin{axis}[bkplot, width=60mm, height=42mm, xmin=0, xmax=2.8, ymin=0, ymax=36, xtick={0.5,1.5,2.5}, xticklabels={$\tau$,$3\tau$,$5\tau$}, ytick={28.5,30}, yticklabels={,30}, xlabel={$t$（ms）}, ylabel={$i$（mA）}]
  \addplot[bkthin, dashed, domain=0:2.6] {30};
  \addplot[bkcurve, domain=0:2.6, samples=60] {30*(1-exp(-x/0.5))};
  \draw[bkGray, dashed, line width=0.5pt] (axis cs:1.5,0) -- (axis cs:1.5,28.5);
  \node[font=\sffamily\footnotesize, anchor=north west, fill=white, inner sep=1pt] at (axis cs:1.55,26) {28.5};
\end{axis}
\end{tikzpicture}
\end{minipage}\hfill
\begin{minipage}[c]{0.46\linewidth}\raggedright
電流は、霧のつぶの速さ（例題5のグラフ）とまったく同じ形でふえ、「電流の終端速度」30\,mAに近づきます。時定数は霧のつぶ（1.2\,ms）とリレー（0.5\,ms）で同じくらいの短さです。
\end{minipage}
\par\vspace{4pt}

\begin{jibun}
$E=\qty{24}{V}$、$R=\qty{60}{\ohm}$、$L=\qty{0.3}{H}$ の RL 回路で、（1）最後の電流、（2）時定数、（3）スイッチを入れて $3\tau$ たったときの電流を求めましょう（$e^{-3}=0.050$）。
\bkjans{（1）\nobreak \qty{0.4}{A}。（2）\nobreak $\dfrac{0.3}{60}=\qty{0.005}{s}=\qty{5}{ms}$。（3）\nobreak $0.4\times0.95=\qty{0.38}{A}$。}
\end{jibun}

\begin{reidai}[力のつり合いで決まる最高速度]
自動車の空気の抵抗は $f=cv^2$、$c=\qty{0.4}{N.s^2/m^2}$ とします。車輪が車を前へおす力が\qty{1000}{N}のとき、空気の抵抗とつり合う速さ（それ以上速くならない速さ）は何m/sですか。時速何kmですか（ほかの抵抗は考えない）。\par
\kangaekata 終端速度と同じ考え。おす力と抵抗がつり合うと、加速度が0になる。重力 $mg$ のかわりに、おす力1000\,Nが入る。\par
\shiki $1000=0.4\,v^2$ より $v^2=\dfrac{1000}{0.4}=2500$、$v=\sqrt{2500}=50$\bkvifq
$\qty{50}{m/s}\times\qty{3600}{s/h}=\qty{180000}{m/h}=\qty{180}{km/h}$\par
\kotae \qty{50}{m/s}（時速180\,km）
\end{reidai}
\noindent\begin{minipage}[c]{0.5\linewidth}
\begin{tikzpicture}
\begin{axis}[bkplot, width=55mm, height=42mm, xmin=0, xmax=64, ymin=0, ymax=1500, xtick={20,40,50,60}, xticklabels={20,40,,60}, ytick={500,1000}, xlabel={$v$（m/s）}, ylabel={力（N）}]
  \addplot[bkcurve, domain=0:61, samples=60] {0.4*x^2};
  \addplot[bkcurve2, domain=0:61] {1000};
  \draw[bkGray, dashed, line width=0.5pt] (axis cs:50,0) -- (axis cs:50,1000);
  \node[font=\sffamily\footnotesize, anchor=north] at (axis cs:50,-80) {50};
\end{axis}
\end{tikzpicture}
\end{minipage}\hfill
\begin{minipage}[c]{0.46\linewidth}\raggedright
青は空気の抵抗 $0.4v^2$、オレンジはおす力1000\,N。50\,m/sより遅いと、おす力のほうが大きいので速くなり、速いと抵抗のほうが大きいので遅くなります。だから、まじわる点の速さに落ち着きます。
\end{minipage}
\par\vspace{4pt}

\begin{jibun}
自転車と人の空気の抵抗を $f=cv^2$、$c=\qty{0.25}{N.s^2/m^2}$ とします。ペダルをこいで前へおす力が\qty{20}{N}のとき、空気の抵抗とつり合う速さは何m/sですか。時速何kmですか（ほかの抵抗は考えない）。
\bkjans{$v=\sqrt{\dfrac{20}{0.25}}=\sqrt{80}\approx\qty{8.9}{m/s}$。$8.94\times3.6\approx32$。時速約32\,km。}
\end{jibun}

%% ---------------- まちがえやすいところ・物理工学 ----------------
\bksection{tsumazuki}{まちがえやすいところをチェック}

\begin{tsumazuki}{上に投げたときの抵抗の向き}
\bkng{上がるときも、下向きを $+$ として $m\dfrac{\dd v}{\dd t}=mg-cv^2$}
\bkok{上がるときは $m\dfrac{\dd v}{\dd t}=mg+cv^2$（抵抗も下向き）}
抵抗の向きは、いつも動きと反対です。上へ動いているときは、重力も抵抗も下向きになり、速さは空気がないときより早くへります。$-cv^2$ は、下へ動くときだけの書き方です。
\end{tsumazuki}

\begin{tsumazuki}{$b$ と $c$ の単位を同じにする}
\bkng{$b$ も $c$ も、単位は N·s/m}
\bkok{$b$ はN·s/m、$c$ はN·s$^2$/m$^2$}
$f=bv$ と $f=cv^2$ で、$f$ はどちらもNです。$v^2$ の単位は $(\text{m/s})^2$ なので、$c$ の単位には s$^2$ と m$^2$ が入ります。単位がちがうので、$b$ と $c$ の大きさをくらべても意味がありません。
\end{tsumazuki}

\begin{tsumazuki}{重い物ほど速く落ちる、はいつも正しい？}
\bkng{重い物は、いつでも速く落ちる}
\bkok{空気がなければ、重さによらず同じ速さで落ちる}
重い物が速く落ちるのは、空気の抵抗があるときだけです。空気のない月の上では、ハンマーと羽根が同時に落ちることが、1971年にアポロ15号の宇宙飛行士の実験で確かめられました（空気がないときの落ち方は5-13）。
\end{tsumazuki}

\begin{tsumazuki}{終端速度に「着く」時刻を求める}
\bkng{$v=v_{\mathrm f}(1-e^{-t/\tau})=v_{\mathrm f}$ となる $t$ を求める}
\bkok{終端速度には限りなく近づくだけ。95\,\%なら $3\tau$、99.3\,\%なら $5\tau$}
$e^{-t/\tau}$ は0になりません（6-4⑤）。「何\,\%になるまで」と決めて、時間を求めます。
\end{tsumazuki}

\begin{tsumazuki}{止まるときは、いつも軽く感じる？}
\bkng{下へ動いているエレベーターが止まるときは、軽く感じる}
\bkok{下へ動いているエレベーターが止まるときは、重く感じる}
下へ動いていて止まるときは、速さ（下向き）がへるので、加速度は上向き（$a>0$）です。$N=m(g+a)$ で $N>mg$ なので、重く感じます。決めるのは、動いている向きではなく、加速度の向きです。
\end{tsumazuki}

\begin{butsuri}
\begin{itemize}[leftmargin=1.2\zw, itemsep=3pt, topsep=0pt, label={\color{bkPurple}●}]
\item \textbf{パラシュート}\par\noindent\begin{minipage}[c]{0.8\linewidth}
パラシュートは、前から見た面積を大きくして $c$ を何十倍にもし、終端速度を小さくする道具です。$v_{\mathrm t}=\sqrt{\dfrac{mg}{c}}$ なので、$c$ を100倍にすると、終端速度は $\frac{1}{10}$ になります。人を乗せたパラシュートは、毎秒5\,mほどで地面に着きます（例題2・じぶんでとこう2）。
\end{minipage}\hfill
\begin{minipage}[c]{0.17\linewidth}\centering
\begin{tikzpicture}[x=1mm, y=1mm]
  \filldraw[fill=bkOrange!45, draw=bkOrange!80!black, line width=0.7pt] (-9,14) arc (180:0:9 and 5.5) -- cycle;
  \foreach \x in {-9,-4.5,0,4.5,9} {\draw[bkInk!50, line width=0.4pt] (\x,14) -- (0,2);}
  \bkperson{(0,-6)}{10}{bkBlue}
\end{tikzpicture}
\end{minipage}
\item \textbf{落ちる玉で、ねばりけをはかる}\par\noindent\begin{minipage}[c]{0.8\linewidth}
ねばりけの強い液体（はちみつや油）に小さな鉄の玉を落とすと、すぐに終端速度になって、一定の速さで沈みます。速さに比例する抵抗の係数 $b$ は液体のねばりけに比例するので、沈む速さをはかれば、ねばりけがわかります。これを落球式の粘度計といい、エンジンオイルなどを調べるのに使われます。
\end{minipage}\hfill
\begin{minipage}[c]{0.17\linewidth}\centering
\begin{tikzpicture}[x=1mm, y=1mm]
  \filldraw[fill=bkYellow!35, draw=bkInk!60, line width=0.6pt] (0,0) rectangle (12,32);
  \fill[bkGray!70] (6,22) circle (1.4);
  \draw[bkInk!60, line width=0.6pt, -{Stealth[length=3pt]}] (6,19) -- (6,13);
  \draw[bkRed, line width=0.5pt] (-1,26) -- (13,26) (-1,8) -- (13,8);
\end{tikzpicture}
\end{minipage}
\end{itemize}
\end{butsuri}

\begin{butsuri}
\begin{itemize}[leftmargin=1.2\zw, itemsep=3pt, topsep=0pt, label={\color{bkPurple}●}]
\item \textbf{自動車の燃費と空気の抵抗}\par\noindent 空気の抵抗は速さの2乗に比例します。$c=\qty{0.4}{N.s^2/m^2}$ の車なら、時速50\,kmで約\qty{77}{N}、時速100\,kmで約\qty{309}{N}（4倍）。速く走ると燃料を多く使うのは、このためです。
\item \textbf{磁石のブレーキ}\par\noindent 磁石のそばを金属の板が速く動くと、板に「うず電流」が流れ、速さに比例する力（$f=bv$ の形）で動きをさまたげます。遊園地のフリーフォールや一部の電車のブレーキに使われ、ふれ合わないのですりへりません。
\end{itemize}
\end{butsuri}

%% ---------------- 章末問題 ----------------
\bkendproblems
\bkproblemgroup{基本}{2}

\bkproblem ある小さなつぶは、$\dfrac{\dd v}{\dd t}=9.8-980\,v$（$v$ はm/s、$t$ は秒）にしたがって落ちます。終端速度と時定数を求めましょう。
\bkpans{$v_{\mathrm f}=\dfrac{9.8}{980}=\qty{0.01}{m/s}$。$\tau=\dfrac{1}{980}\approx\qty{0.0010}{s}$。}

\bkproblem 卓球のボール（質量0.0027\,kg）の空気の抵抗は $f=cv^2$、$c=0.00033$（N·s$^2$/m$^2$）です。終端速度は何m/sですか。
\bkpans{$\sqrt{\dfrac{0.0027\times9.8}{0.00033}}=\sqrt{80.2}\approx\qty{9.0}{m/s}$。}

\bkproblem 形と大きさが同じで、2乗に比例する抵抗を受けて落ちる物の質量を25倍にすると、終端速度は何倍になりますか。
\bkpans{$\sqrt{25}=5$ 倍。}

\bkproblem $\dfrac{\dd v}{\dd t}=9.8-0.004\,v^2$、はじめ $v=0$ として、$\Delta t=1$ 秒のオイラー法で、1秒後と2秒後の速さを求めましょう。
\bkpans{1秒後 $\qty{9.8}{m/s}$。加速度 $9.8-0.004\times96.04\approx9.416$、2秒後 $9.8+9.416\approx\qty{19.2}{m/s}$。}

\bkproblem 体重70\,kgの人が、エレベーターの中で体重計に乗っています。加速度が（1）上向きに $\qty{0.5}{m/s^2}$、（2）下向きに $\qty{0.5}{m/s^2}$ のとき、目もりは何kgですか（$g=9.8$）。
\bkpans{（1）\nobreak $70\times\dfrac{10.3}{9.8}\approx73.6$\,kg。（2）\nobreak $70\times\dfrac{9.3}{9.8}\approx66.4$\,kg。}

\bkproblem $E=\qty{5}{V}$、$R=\qty{25}{\ohm}$、$L=\qty{0.05}{H}$ の RL 回路で、（1）最後の電流、（2）時定数、（3）スイッチを入れて $3\tau$ たったときの電流を求めましょう（$e^{-3}=0.050$）。
\bkpans{（1）\nobreak \qty{0.2}{A}。（2）\nobreak $\dfrac{0.05}{25}=\qty{0.002}{s}=\qty{2}{ms}$。（3）\nobreak $0.2\times0.95=\qty{0.19}{A}$。}

\bkproblemgroup{チャレンジ}{3}

\bkproblem エレベーターの中で体重計に乗りました。目もりから、加速度の向きと大きさを求めましょう（$g=9.8$）。\newline（1）体重80\,kgの人で、目もりが84.0\,kg。\newline（2）体重60\,kgの人で、目もりが57.0\,kg。
\bkpans{（1）\nobreak $\dfrac{84}{80}=1.05=\dfrac{9.8+a}{9.8}$ より $a\approx0.49$。上向きに約\qty{0.49}{m/s^2}。（2）\nobreak $\dfrac{57}{60}=0.95=\dfrac{9.8+a}{9.8}$ より $a\approx-0.49$。下向きに約\qty{0.49}{m/s^2}。}

\bkproblem 問題1のつぶが、止まった状態から落ち始めます。$v$ を $t$ の式で表し、$3\tau$ たったときの速さを求めましょう（$e^{-3}=0.050$）。
\bkpans{$v=0.01\,(1-e^{-980t})$。$3\tau$ で $0.01\times0.95=\qty{0.0095}{m/s}$。}

\bkproblem 空気の抵抗が $f=cv^2$、$c=\qty{0.35}{N.s^2/m^2}$ の自動車を\qty{700}{N}でおすとき、抵抗とつり合う速さは何m/s、時速何kmですか（ほかの抵抗は考えない）。
\bkpans{$v=\sqrt{\dfrac{700}{0.35}}=\sqrt{2000}\approx\qty{44.7}{m/s}$。$44.7\times3.6\approx161$。時速約161\,km。}

\bkproblem 空気の抵抗が $f=cv^2$ のとき、終端速度は何m/sですか。\newline（1）質量\qty{100}{kg}のスカイダイバーが、手足を広げて落ちているとき。$c=\qty{0.25}{N.s^2/m^2}$ とし、時速何kmかも答えましょう。\newline（2）質量\qty{70}{kg}の人が、パラシュートで降りているとき（$c=\qty{25}{N.s^2/m^2}$）。
\bkpans{（1）\nobreak $v_{\mathrm t}=\sqrt{\dfrac{100\times9.8}{0.25}}=\sqrt{3920}\approx\qty{62.6}{m/s}$。$62.6\times3.6\approx225$。時速約225\,km。（2）\nobreak $v_{\mathrm t}=\sqrt{\dfrac{70\times9.8}{25}}=\sqrt{27.44}\approx\qty{5.2}{m/s}$。}

\begin{furikaeri}
\begin{checklist}
\item 落ちる物の式を立て、終端速度を求められる。\bkfurref{bk:6-5:shuutan}
\item 表計算で、落ちる物の速さを計算できる。\bkfurref{bk:6-5:hyou}
\item 速さに比例する抵抗の式を解ける（RL 回路も）。\bkfurref{bk:6-5:toku}
\end{checklist}
\end{furikaeri}

\bkfukidashi{mina}{雨つぶがやさしく落ちてくるのは、空気の抵抗のおかげだった！　次は、ばねにつるしたおもりの「ゆれ」を式で調べよう。}
